\documentclass[10pt,a4paper]{amsart}
\usepackage[margin=19mm]{geometry}
\usepackage{amsmath,amssymb,mathtools}
\usepackage[scr=rsfs,cal=euler]{mathalfa}
\usepackage{xfrac}
\usepackage[unicode,hidelinks,bookmarks=false]{hyperref}
\newcommand{\ie}{i.e.\ }
\newcommand{\cf}{cf.\ }
\newcommand{\resp}{respectively\ }
\newcommand{\Subsection}{\subsection}
\providecommand{\nobreakdash}{\nobreak\hbox{-}}
\setlength{\emergencystretch}{2em}
\multlinegap0pt
\usepackage{amssymb}
\newtheorem{proposition}{Proposition}
\title{Bifurcations of solitary waves \\ in a coupled system of long and short waves}
\author{James Hornick, Dmitry E. Pelinovsky}
\date{Source: arXiv:2603.28531v1 (CC BY 4.0)}
\begin{document}
\maketitle

\noindent\emph{Source and licence.} Bifurcations of solitary waves \\ in a coupled system of long and short waves. Version arXiv:2603.28531v1 (CC BY 4.0), distributed by arXiv under the Creative Commons Attribution 4.0 International licence, http://creativecommons.org/licenses/by/4.0/, as recorded in the arXiv licence metadata for this exact version and shown on the abstract page https://arxiv.org/abs/2603.28531v1. Full licence notice: \texttt{licenses/arxiv-2603.28531-CC-BY-4.0.txt}.\par\smallskip

\noindent\textit{Editorial scope.} Counted unit: the article's proposition theorem-Morse-first (source0.tex lines 687--692). The article's complete original proof of it is delivered with it (source0.tex lines 694--817, one \texttt{proof} environment of 124 lines). No mathematics is written, repaired or re-derived: the statement and the proof are the article's own lines, in the article's own notation, and no retained line is substituted or re-flowed.  Retained uncounted as the source-local context the proof needs: lines 202--205; lines 385--391; lines 433--443; lines 507--513; lines 558--561; lines 577--588; lines 662--665; lines 667--670; lines 674--677.  Nothing was omitted from any retained span, so the package carries no omission record.  No retained line contains a citation, so the wrapper carries no bibliography and no bibliography entry.  No retained line contains a figure environment, an image-inclusion command or drawing code, so no image is delivered and the package declares no asset dependency.  Editorial formatting only, in addition: an AMS \texttt{amsart} preamble that restates the article's own declarations and macros used by the retained lines, namely the environments \texttt{proposition} on separate counters, together with the standard packages those lines need.  The retained environments print in this document as Proposition 1 in order of appearance, whereas the article prints the same items with the numbering of its own full document.  The complete native source of the article as distributed in the arXiv e-print for this exact version is bundled unchanged as \texttt{sources/197452/source0.tex}.
\begin{equation}
\label{solution-1}
U = 3 c \; \text{sech}^2\left(\frac{\sqrt{c}}{2}\xi\right), \quad A=0,
\end{equation}

\begin{equation}
    \label{Hessian}
\mathcal{L}=-\begin{pmatrix} \partial_{\xi}^2 + U - c & 2 s A & 0 \\ 
2 s A & \frac{2 s}{k}\left(\partial_{\xi}^2+kU+\Omega\right) & 0 \\
0&  0  &  \frac{2 s}{k}\left(\partial_{\xi}^2+kU+\Omega\right)
\end{pmatrix}.
\end{equation}

\begin{equation}
    \label{block-1}
    \mathcal{L}_J = 
\begin{pmatrix} L_1 & -2 s A  \\ 
-2 s A & L_2 
\end{pmatrix} \qquad
\mbox{\rm with} \;\; 
\left\{ \begin{array}{l} L_1 = -\partial_{\xi}^2 + c - U, \\
L_2 = \frac{2 s}{k} \left( - \partial_{\xi}^2 - \Omega - kU \right),
\end{array} \right.
\end{equation}

\begin{align}
\label{block-2}
\left\{ \begin{array}{l}
    L_1 = \frac{c}{4} \left( - \partial_y^2 + 4 - 12 {\rm sech}^2(y) \right), \\
    L_2 = \frac{c}{2 |k|} \left( - \partial_y^2 + \frac{4 |\Omega|}{c} - 12 k {\rm sech}^2(y) \right).
    \end{array} \right.
\end{align}

\begin{equation}
\label{g-eigenfunction}
    g(\xi) = \text{sech}^{p}\left( \frac{\sqrt{c}}{2} \xi\right), \quad p = \frac{\sqrt{1+48 k}-1}{2}.
\end{equation}

\begin{proposition}
    \label{theorem-bifurcation-first}
    Assume $c > 0$, $\Omega < 0$, and $s = {\rm sgn}(k) = 1$.
    Let $U_0$ be given by (\ref{solution-1}) and $g$ be given by (\ref{g-eigenfunction}). 
    Assume that $\langle g^2, L_1^{-1} g^2 \rangle \neq 0$, where $L_1$ is defined in (\ref{block-2}). 
    For every $\Omega$ near $\Omega_c$ such that ${\rm sgn}(\Omega - \Omega_c) = -{\rm sgn}(\langle g^2, L_1^{-1} g^2 \rangle)$, 
    there exists a unique family of solutions with the profile $(U,A) \in H^2_{\rm even}(\mathbb{R}) \times H^2_{\rm even}(\mathbb{R})$ satisfying 
    $$
    \| U - U_0 \|_{H^2} \leq C |\Omega - \Omega_c|, \quad \| A \| \leq C \sqrt{|\Omega - \Omega_c|},
    $$
    for some $\Omega$-independent constant $C > 0$.
\end{proposition}

\begin{equation}
    \label{bounds-on-terms}
\| \mathfrak{w}_1(a) \|_{H^2} \leq C a^2, \qquad \| \mathfrak{z}_1(a) \|_{H^2} \leq C |a|^3,
\end{equation}

\begin{equation}
    \label{near-identity}
\mathfrak{w}_1(a) = a^2 w_2 + \mathfrak{w}_2(a), \quad \| \mathfrak{w}_2(a) \|_{H^2} \leq C a^4, 
\end{equation}

\begin{equation}
    \label{delta-Omega}
\delta \Omega = - k a^2  \frac{\langle g^2, L_1^{-1} g^2 \rangle}{\| g \|^2_{L^2}} + \mathcal{O}(a^4),
\end{equation}

\begin{proposition}
    \label{theorem-Morse-first}
    Let $(U,A)$ be the profile of the bifurcating branch for $\Omega$ near $\Omega_c$ given by Proposition \ref{theorem-bifurcation-first}. 
Its Morse index is equal to $1$ if $\langle g^2, L_1^{-1} g^2 \rangle < 0$ and to $2$ if $\langle g^2, L_1^{-1} g^2 \rangle > 0$, 
    whereas the nullity index is equal to $2$ in both cases. 
\end{proposition}

\begin{proof}
We use the construction of Proposition \ref{theorem-bifurcation-first} with the profile 
\begin{equation}
\label{decomposition-bif}
U = U_0 + \mathfrak{w}_1(a) \quad  \mbox{\rm and} \quad 
A = a g + \mathfrak{z}_1(a),
\end{equation}
where $(\mathfrak{w}_1(a), \mathfrak{z}_1(a)) \in H^2_{\rm even}(\mathbb{R}) \times H^2_{\rm even}(\mathbb{R})$ satisfy the bounds (\ref{bounds-on-terms}). The parameter $a > 0$ parameterize the bifurcating branch, 
in particular, we have $\Omega = \Omega_c + \delta \Omega(a)$, where $\delta \Omega(a)$ is given by the expansion 
(\ref{delta-Omega}) with ${\rm sgn}(\delta \Omega) = -{\rm sgn}(\langle g^2, L_1^{-1} g^2 \rangle)$. 
The Hessian block $L_J$ defined by (\ref{block-1}) can be expanded by using (\ref{decomposition-bif}) as follows:
\begin{align*}
    L_J &= \begin{pmatrix} L_1 - \mathfrak{w}_1(a) & -2 s (a g + \mathfrak{z}_1(a))  \\ 
-2 s (a g + \mathfrak{z}_1(a)) & L_2 - 2 s \mathfrak{w}_1(a) -\frac{2 s}{k} \delta\Omega(a)
\end{pmatrix} \\
&= \begin{pmatrix} L_1 & 0  \\ 
0 & L_2 \end{pmatrix} 
+ a \begin{pmatrix} 0 & -2 s g \\ 
-2 s g & 0 \end{pmatrix} 
+ a^2 \begin{pmatrix} - s L_{1}^{-1} g^2 & 0 \\ 
0 & -2 L_{1}^{-1}g^2+  2\frac{\langle g^2, L_1^{-1} g^2 \rangle}{\| g \|^2_{L^2}}  
\end{pmatrix} + \mathcal{O}(a^3) \\
&= L_J^{(0)} + a L_J^{(1)} + a^2 L_J^{(2)} + \mathcal{O}(a^3),
\end{align*}
where we have used the leading-order terms for $\mathfrak{w}_1(a)$ and $\delta \Omega(a)$ from (\ref{near-identity}) and (\ref{delta-Omega}).
We are now looking at the eigenvalues and eigenvectors of the eigenvalue equation $L_J v = \lambda v$. 
Since $0$ is a double eigenvalue of $L_J^{(0)}$, we are looking for the splitting of the zero eigenvalue 
by using the perturbation theory. We write 
$$
v = v^{(0)} + a v^{(1)} + a^2 v^{(2)} + \mathcal{O}(a^3), \quad \lambda = a \lambda^{(1)} + a^2 \lambda^{(2)} + \mathcal{O}(a^3),
$$
where
$$
v^{(0)} = c_1 \begin{pmatrix} U_0'\\0\end{pmatrix} + c_2\begin{pmatrix}0\\g\end{pmatrix}
$$
for some $(c_1,c_2) \in \mathbb{R}^2$. 


\underline{At the order of $\mathcal{O}(a)$,} we obtain 
\[
L^{(0)} v^{(1)} + L^{(1)} v^{(0)} =\lambda^{(1)} v^{(0)}.
\]
Writing in the component form, we obtain 
$$
v_1=\begin{pmatrix}f_1\\g_1\end{pmatrix} : \quad \left\{ \begin{array}{l} L_1 f_1 - 2 s g (c_2 g) = \lambda^{(1)} (c_1 U_0'), \\
L_2 g_1 - 2 s g (c_1 U_0') = \lambda^{(1)} (c_2 g). \end{array} \right.
$$
Since even $g^2$ is orthogonal to $U_0'$ and odd $g U_0'$ is orthogonal to even $g$, we obtain by Fredholm's theory 
for linear inhomogeneous equations that $\lambda^{(1)} = 0$. This yields the exact solution in the form 
$$
f_1 = 2 s c_2 L_1^{-1} g^2, \quad g_1 = 2 s c_1 L_2^{-1} g U_0',
$$
where $L_1^{-1}$ is uniquely defined in the space of even functions and $L_2^{-1}$ is uniquely defined in the space of odd functions. 
Differentiating of $(-\partial_{\xi}^2 - \Omega_c - k U_0) g = 0$ in $\xi$ yields $L_2 g' = 2 s U_0' g$, which ensures that 
$g_1 = c_1 g'$.

\underline{At the order of $\mathcal{O}(a^2)$,} we obtain 
\[
L^{(0)} v^{(2)} + L^{(1)} v^{(1)} + L^{(2)} v^{(0)} =\lambda^{(2)} v^{(0)}.
\]
Writing in the component form, we obtain 
$$
v_2=\begin{pmatrix}f_2\\g_2\end{pmatrix} : \quad \left\{ \begin{array}{l} L_1 f_2 - 2 s g g_1 - s  (L_1^{-1} g^2) (c_1 U_0') = \lambda^{(2)} (c_1 U_0'), \\
L_2 g_2 - 2 s g f_1 + \left( - 2 L_1^{-1} g^2 + 2 \frac{\langle g^2, L_1^{-1} g^2 \rangle}{\| g \|^2_{L^2}} \right) (c_2 g) = \lambda^{(2)} (c_2 g). \end{array} \right.
$$
Taking the inner product of the first equation with \(U_0'\in \text{Ker}(L_1)\) and 
also taking the inner product of the second equation with \(g\in \text{Ker}(L_2)\), we get the linear system 
for $(c_1,c_2) \in \mathbb{R}^2$:
\begin{align*}
-4 c_1 \langle g U_0', L_2^{-1} g U_0' \rangle - s c_1 \langle (U'_0)^2, L_{1}^{-1}g^2 \rangle &= \lambda^{(2)} c_1 \Vert U_{0}'\Vert^2_{L^2}, \\
-4c_2\langle g^2, L_{1}^{-1}g^2\rangle &= \lambda^{(2)} c_2 \Vert g\Vert^2_{L^2}
\end{align*}
where we have used the explicit expressions for $f_1$, $g_1$. Since the linear system is diagonal, we get two different solutions 
related to the subspaces $(c_1,c_2) = (1,0)$ and $(c_1,c_2) = (0,1)$. 

For the subspace $(c_1,c_2) = (1,0)$, we prove that 
\begin{equation}
    \label{constraint-transl}
4 \langle g U_0', L_2^{-1} g U_0' \rangle + s \langle (U'_0)^2, L_{1}^{-1}g^2 \rangle = 0,
\end{equation}
which yields $\lambda^{(2)} = 0$. This is consistent with the fact that the zero eigenvalue $\lambda = 0$ is 
preserved along the solution branch with the profile $(U,A)$ by the translational symmetry. To verify (\ref{constraint-transl}), we derive from 
\begin{align*}
    U_0'' &= c U_0 - \frac{1}{2} U_0^2, \\
    (U_0')^2 &= c U_0^2 - \frac{1}{3} U_0^3,
\end{align*}
that 
\begin{align*}
    L_1 U_0 &= -U_0'' + c U_0 - U_0^2 = -\frac{1}{2} U_0^2, \\
    L_1 U_0^2 &= -2 U_0 U_0'' - 2 (U_0')^2 + c U_0^2 - U_0^3 = -3 c U_0^2 + \frac{2}{3} U_0^3.
\end{align*}
Hence we obtain 
\begin{align*}
\langle (U'_0)^2, L_{1}^{-1}g^2 \rangle &= c \langle U_0^2, L_{1}^{-1}g^2 \rangle - \frac{1}{3} \langle U_0^3, L_{1}^{-1}g^2 \rangle\\
 &= c \langle L_{1}^{-1} U_0^2, g^2 \rangle - \frac{1}{3} \langle L_{1}^{-1} U_0^3, g^2 \rangle\\
  &= -\frac{c}{2} \langle L_{1}^{-1} U_0^2, g^2 \rangle - \frac{1}{2} \langle U_0^2, g^2 \rangle\\
    &= c \langle U_0, g^2 \rangle - \frac{1}{2} \langle U_0^2, g^2 \rangle.
\end{align*}
On the other hand, since $L_2^{-1} g U_0' = \frac{s}{2} g'$, we also obtain 
\begin{align*}
4 \langle g U_0', L_2^{-1} g U_0' \rangle &= 2 s \langle g U_0', g' \rangle \\
 &= - s \langle U_0'', g^2 \rangle \\
  &= -c s  \langle U_0, g^2 \rangle + \frac{s}{2} \langle U_0^2, g^2 \rangle.
\end{align*}
Substituting these computations into the left-hand side of (\ref{constraint-transl}), we obtain the zero result.

For the subspace $(c_1,c_2) = (0,1)$, we get 
\[
\lambda^{(2)}=-4\frac{\langle g^2,L_{1}^{-1}g^2\rangle}{\Vert g\Vert^2_{L^2}}
\]
which shows that $\lambda = a^2 \lambda^{(2)} + \mathcal{O}(a^3) > 0$ if $\langle g^2,L_{1}^{-1}g^2\rangle < 0$ and $\lambda = a^2 \lambda^{(2)} + \mathcal{O}(a^3) < 0$ if $\langle g^2,L_{1}^{-1}g^2\rangle > 0$. Since $L_J^{(0)}$ has only one negative eigenvalue at $U = U_0$, the splitting of the double zero eigenvalue 
of $L_J^{(0)}$ is clarified above, and the other eigenvalues of $L_J^{(0)}$ are strictly positive, the continuity of eigenvalues in $a$ implies that the Morse index of $L_J$ is $1$ if $\langle g^2,L_{1}^{-1}g^2\rangle < 0$ and $2$ if $\langle g^2,L_{1}^{-1}g^2\rangle > 0$ for small $a \neq 0$. 

Referring back to the Hessian operator $\mathcal{L}$ given by (\ref{Hessian}), it follows that 
the operator 
$$
L_2 = -\frac{2s}{k} (\partial_{\xi}^2 + k U + \Omega)
$$ 
admits a simple zero eigenvalue due to the rotational symmetry with $L_2 A = 0$. The rest of $L_2$ at $U = U_0$ and $\Omega = \Omega_c$ is strictly positive, 
hence the Morse index of $L_2$ is $0$ for small $a$. This yields the claim on the Morse index of $\mathcal{L}$.
On the other hand, $\mathcal{L}$ has a double zero eigenvalue associated with the translational and rotational symmetries, and the splitting 
of the double zero eigenvalue of $L_J^{(0)}$ clarified above shows that the kernel of $\mathcal{L}$ is exactly double for small $a$. This yields the 
claim about the nullity index of $\mathcal{L}$.
\end{proof}

\end{document}
